\documentclass{article}

\usepackage[utf8]{inputenc}
\usepackage[portuguese]{babel}
\usepackage[T1]{fontenc}
\usepackage{lmodern}
\usepackage{fullpage}
\usepackage[usenames]{color}
\usepackage[color]{coqdoc}
\usepackage{url}
\usepackage{makeidx,hyperref}
\usepackage[all]{xy}

\usepackage{mathpartir}
\usepackage{tcolorbox}
\usepackage{amsmath,amssymb}

\newcommand{\tto}{\twoheadrightarrow}
\newcommand{\ott}{\twoheadleftarrow}

\title{Formalização da correção do algoritmo Selection Sort}
\author{Nome e matrícula}
\date{\today}

\begin{document}
\maketitle
\section{Dados dos alunos}
\begin{itemize}
      \item Luís Eduardo Curi Serra  - 251033574
      \item Weldo Gonçalves da Silva Junior - 222014133
\end{itemize}
\section{Introdução}

Este relatório documenta a formalização, no assistente de provas Rocq (anteriormente
conhecido como Coq), da correção do algoritmo de ordenação por seleção
(\emph{Selection Sort}), proposta como Tema 2 do Projeto de Formalização da
disciplina de Projeto e Análise de Algoritmos (2026/1).

A correção de um algoritmo de ordenação é caracterizada por duas propriedades
que devem valer simultaneamente sobre a lista de saída:

\begin{itemize}
  \item \textbf{Ordenação}: a lista resultante está ordenada (\texttt{Sorted le}).
  \item \textbf{Permutação}: a lista resultante contém exatamente os mesmos
        elementos da lista de entrada, com as mesmas multiplicidades
        (\texttt{Permutation}).
\end{itemize}

O teorema final demonstrado é:
\[
\forall l,\; \mathtt{Sorted\ le\ (ss\ l)} \land \mathtt{Permutation\ l\ (ss\ l)}
\]

\section{O algoritmo Selection Sort}

A ideia do algoritmo é simples: a cada etapa, encontra-se o menor elemento
ainda não posicionado, e ele é movido para a frente da lista resultante. O
processo se repete recursivamente sobre o restante da lista.

A função auxiliar \texttt{select\_min} percorre a lista comparando elementos
consecutivos e retorna o menor elemento encontrado (ou \texttt{None} se a
lista for vazia). Como a recursão de \texttt{select\_min} não é estrutural —
a chamada recursiva é feita sobre uma lista reconstruída (\texttt{h1::tl} ou
\texttt{h2::tl}), e não sobre uma subestrutura direta do argumento original —
foi necessário utilizar o comando \texttt{Function} com uma medida
(\texttt{measure length l}) explícita, provando que essa medida decresce a
cada chamada recursiva.

\section{Erro identificado no enunciado}
\label{sec:bug}

Conforme incentivado pelo enunciado do projeto, documentamos aqui um erro
encontrado na definição original de \texttt{ss} proposta no enunciado:

\begin{verbatim}
Fixpoint ss (l: list nat) :=
  match l with
  | nil => nil
  | h::tl => match select_min l with
             | None => nil
             | Some m => m::(ss tl)
             end
  end.
\end{verbatim}

O problema está na chamada recursiva \texttt{ss tl}: ela recursiona sobre a
\emph{cauda} da lista (\texttt{tl}), e não sobre a lista da qual o mínimo
\texttt{m} foi de fato removido. Isso faz com que o elemento mínimo
encontrado (\texttt{m}) permaneça na lista usada na próxima chamada
recursiva, caso ele não seja a própria cabeça \texttt{h}.

Um contraexemplo concreto evidencia o problema. Considere
\texttt{l = [4;2;1;3]}:

\begin{itemize}
  \item \texttt{select\_min [4;2;1;3] = Some 1}, então o resultado começa com
        \texttt{1 ::}.
  \item A recursão, no entanto, é chamada sobre \texttt{tl = [2;1;3]} (a
        cauda de \texttt{l}), e \emph{não} sobre \texttt{[4;2;3]} (a lista
        sem o mínimo \texttt{1}).
  \item Isso repete o elemento \texttt{1} nas chamadas seguintes e descarta
        o elemento \texttt{4} definitivamente, produzindo
        \texttt{ss [4;2;1;3] = [1;1;1;3]}.
\end{itemize}

O resultado \texttt{[1;1;1;3]} não é uma permutação de \texttt{[4;2;1;3]}
(o elemento \texttt{4} desaparece e o \texttt{1} é triplicado). Logo, o
teorema \texttt{selectionsort\_correct}, como enunciado sobre esta definição
de \texttt{ss}, é \textbf{falso} — não haveria como demonstrá-lo em Rocq.

\subsection{Correção proposta}

A correção consiste em fazer a recursão sobre a lista da qual o mínimo foi
efetivamente removido. Para isso, introduzimos a função auxiliar
\texttt{remove\_one}, que remove uma única ocorrência de um elemento de uma
lista:

\begin{verbatim}
Fixpoint remove_one (x : nat) (l : list nat) : list nat :=
  match l with
  | nil => nil
  | h :: tl => if h =? x then tl else h :: remove_one x tl
  end.
\end{verbatim}

E redefinimos \texttt{ss} como:

\begin{verbatim}
Function ss (l : list nat) {measure length l} : list nat :=
  match select_min l with
  | None => nil
  | Some m => m :: ss (remove_one m l)
  end.
\end{verbatim}

Com esta correção, \texttt{ss [4;2;1;3]} computa corretamente para
\texttt{[1;2;3;4]} (verificado por \texttt{Eval compute}, Seção
\ref{sec:testes}).

\section{Estratégia de formalização}

A prova foi organizada em quatro blocos de lemas auxiliares, culminando no
teorema principal.

\subsection{Corretude de \texttt{select\_min}}

Dois lemas caracterizam o comportamento de \texttt{select\_min}:

\begin{itemize}
  \item \texttt{select\_min\_correct}: se \texttt{select\_min l = Some m},
        então \texttt{m} é menor ou igual a todos os elementos de \texttt{l}
        (formalizado pelo predicado auxiliar \texttt{le\_all}).
  \item \texttt{select\_min\_in}: se \texttt{select\_min l = Some m}, então
        \texttt{m} pertence a \texttt{l}. Este lema é necessário para
        justificar que \texttt{remove\_one m l} de fato diminui o tamanho da
        lista, garantindo a terminação de \texttt{ss}.
  \item \texttt{select\_min\_none}: se \texttt{select\_min l = None}, então
        \texttt{l} é vazia — usado no caso base da prova do teorema final.
\end{itemize}

Ambas as provas são feitas por \emph{indução funcional}
(\texttt{functional induction}), tática que utiliza o princípio de indução
derivado automaticamente da definição recursiva de \texttt{select\_min},
gerando exatamente os casos correspondentes aos ramos do \texttt{match}/
\texttt{if} da definição original.

\subsection{Propriedades de \texttt{remove\_one}}

Três lemas caracterizam a função de remoção:

\begin{itemize}
  \item \texttt{remove\_one\_length}: remover um elemento presente na lista
        diminui estritamente o seu tamanho — essencial para justificar a
        medida de terminação de \texttt{ss}.
  \item \texttt{remove\_one\_In}: todo elemento presente após a remoção já
        estava presente na lista original.
  \item \texttt{remove\_one\_perm}: se \texttt{x} pertence a \texttt{l},
        então \texttt{l} é uma permutação de \texttt{x :: remove\_one x l}.
        A prova utiliza o construtor \texttt{perm\_swap} de
        \texttt{Permutation} para trazer o elemento removido do meio da
        lista para a posição inicial.
\end{itemize}

Todas as três provas são feitas por indução estrutural simples sobre a
lista \texttt{l}.

\subsection{Lema auxiliar de ordenação}

O lema \texttt{le\_all\_sorted} estabelece que, se \texttt{l} está ordenada
e \texttt{x} é menor ou igual a todos os elementos de \texttt{l}, então
\texttt{x :: l} também está ordenada. Este lema é o elo entre a corretude de
\texttt{select\_min} (que garante \texttt{le\_all}) e a propriedade
\texttt{Sorted} exigida pelo teorema final.

\subsection{Teorema principal}

A prova de \texttt{selectionsort\_correct} também é feita por indução
funcional, agora sobre a definição de \texttt{ss}, gerando dois casos:

\begin{itemize}
  \item \textbf{Caso \texttt{select\_min l = None}}: pelo lema
        \texttt{select\_min\_none}, conclui-se que \texttt{l = nil}, e as
        duas propriedades (\texttt{Sorted} e \texttt{Permutation}) valem
        trivialmente para a lista vazia.
  \item \textbf{Caso \texttt{select\_min l = Some m}}: a hipótese de indução
        garante que \texttt{ss (remove\_one m l)} já está ordenada e é uma
        permutação de \texttt{remove\_one m l}. A partir daí:
        \begin{itemize}
          \item A \textbf{permutação} de \texttt{l} para
                \texttt{m :: ss (remove\_one m l)} é obtida encadeando
                \texttt{remove\_one\_perm} com a hipótese de indução, via
                transitividade de \texttt{Permutation}.
          \item A \textbf{ordenação} é obtida aplicando
                \texttt{le\_all\_sorted}, para o que é necessário mostrar
                que \texttt{m} é menor ou igual a todos os elementos de
                \texttt{ss (remove\_one m l)}. Isso é feito "atravessando"
                a permutação dada pela hipótese de indução (com o lema de
                biblioteca \texttt{Permutation\_in}) para reduzir o problema
                a \texttt{le\_all m l}, já estabelecido por
                \texttt{select\_min\_correct}.
        \end{itemize}
\end{itemize}

\section{Testes e verificação}
\label{sec:testes}

Além da prova formal, verificamos o comportamento computacional de
\texttt{ss} sobre exemplos concretos, tanto por avaliação direta quanto por
lemas demonstrados por \texttt{reflexivity} (o que garante que o resultado
computado bate exatamente com o esperado, e não apenas "parece certo" ao
olho):

\begin{verbatim}
Eval compute in ss [4;2;1;3].
     = [1; 2; 3; 4]

Eval compute in ss [5;4;3;2;1].
     = [1; 2; 3; 4; 5]

Eval compute in ss [].
     = []

Lemma teste1 : ss [4;2;1;3] = [1;2;3;4].
Proof. reflexivity. Qed.

Lemma teste2 : ss [5;4;3;2;1] = [1;2;3;4;5].
Proof. reflexivity. Qed.

Lemma teste3 : ss (@nil nat) = nil.
Proof. reflexivity. Qed.
\end{verbatim}

Adicionalmente, executamos o comando \texttt{Print Assumptions
selectionsort\_correct.} para confirmar que a prova do teorema principal não
depende de nenhum \texttt{Admitted} ou axioma além dos já utilizados
pela biblioteca padrão do Rocq.

\section{Conclusão}

A formalização apresentada corrige um erro presente na definição original
de \texttt{ss} proposta no enunciado (Seção~\ref{sec:bug}), introduzindo a
função auxiliar \texttt{remove\_one} para garantir que o elemento mínimo
seja de fato removido da lista antes da chamada recursiva. Com essa
correção, o teorema de corretude do Selection Sort
(\texttt{selectionsort\_correct}) foi demonstrado por completo em Rocq,
sem lacunas, estabelecendo tanto a propriedade de ordenação quanto a de
permutação da lista de saída em relação à lista de entrada.

\section{Código fonte completo}

A seguir, o arquivo \texttt{selection\_sort.v} completo, com comentários
formatados via \texttt{coqdoc}.

\begin{coqdoccode}
\coqdocnoindent
\coqdockw{From} \coqdocvar{Stdlib} \coqdockw{Require} \coqdockw{Import} \coqdocvar{Arith} \coqdocvar{List} \coqdocvar{Lia}.\coqdoceol
\coqdocnoindent
\coqdockw{From} \coqdocvar{Stdlib} \coqdockw{Require} \coqdockw{Import} \coqdocvar{Recdef}.\coqdoceol
\coqdocnoindent
\coqdockw{From} \coqdocvar{Stdlib} \coqdockw{Require} \coqdockw{Import} \coqdocvar{Sorted}.\coqdoceol
\coqdocnoindent
\coqdockw{From} \coqdocvar{Stdlib} \coqdockw{Require} \coqdockw{Import} \coqdocvar{Permutation}.\coqdoceol
\coqdocnoindent
\coqdockw{Import} \coqdocvar{ListNotations}.\coqdoceol
\coqdocemptyline
\end{coqdoccode}
\coqdoclemma{select\_min} recebe uma lista de naturais e retorna o menor elemento
    desta lista. Se a lista for vazia, \coqdoclemma{select\_min} \coqexternalref{nil}{http://rocq-prover.org/doc/V9.2/corelib//Corelib.Init.Datatypes}{\coqdocconstructor{nil}} retorna None. 
\begin{coqdoccode}
\coqdocemptyline
\coqdocnoindent
\coqdockw{Function} \coqdocvar{select\_min} (\coqdocvar{l} : \coqdocvar{list} \coqdocvar{nat}) \{\coqdockw{measure} \coqdocvar{length} \coqdocvar{l}\} : \coqdocvar{option} \coqdocvar{nat} :=\coqdoceol
\coqdocindent{1.00em}
\coqdockw{match} \coqdocvar{l} \coqdockw{with}\coqdoceol
\coqdocindent{1.00em}
\ensuremath{|} \coqdocvar{nil} \ensuremath{\Rightarrow} \coqdocvar{None}\coqdoceol
\coqdocindent{1.00em}
\ensuremath{|} \coqdocvar{h}::\coqdocvar{nil} \ensuremath{\Rightarrow} \coqdocvar{Some} \coqdocvar{h}\coqdoceol
\coqdocindent{1.00em}
\ensuremath{|} \coqdocvar{h1}::\coqdocvar{h2}::\coqdocvar{tl} \ensuremath{\Rightarrow} \coqdockw{if} \coqdocvar{h1} <=? \coqdocvar{h2} \coqdockw{then} \coqdocvar{select\_min} (\coqdocvar{h1}::\coqdocvar{tl}) \coqdockw{else} \coqdocvar{select\_min} (\coqdocvar{h2}::\coqdocvar{tl})\coqdoceol
\coqdocindent{1.00em}
\coqdockw{end}.\coqdoceol
\coqdocemptyline
\coqdocnoindent
\coqdockw{Definition} \coqdocvar{le\_all} \coqdocvar{x} \coqdocvar{l} := \coqdockw{\ensuremath{\forall}} \coqdocvar{y}, \coqdocvar{In} \coqdocvar{y} \coqdocvar{l} \ensuremath{\rightarrow} \coqdocvar{x} \ensuremath{\le} \coqdocvar{y}.\coqdoceol
\coqdocemptyline
\end{coqdoccode}
Se \coqdoclemma{select\_min} \coqdocvar{l} retorna um natural \coqdocvar{m}, então \coqdocvar{m} é menor ou igual a
    todos os elementos de \coqdocvar{l}. 
\begin{coqdoccode}
\coqdocnoindent
\coqdockw{Lemma} \coqdocvar{select\_min\_correct} : \coqdockw{\ensuremath{\forall}} \coqdocvar{l} \coqdocvar{m}, \coqdocvar{select\_min} \coqdocvar{l} = \coqdocvar{Some} \coqdocvar{m} \ensuremath{\rightarrow} \coqdocvar{le\_all} \coqdocvar{m} \coqdocvar{l}.\coqdoceol
\coqdocemptyline
\end{coqdoccode}
Se \coqdoclemma{select\_min} \coqdocvar{l} retorna \coqexternalref{Some}{http://rocq-prover.org/doc/V9.2/corelib//Corelib.Init.Datatypes}{\coqdocconstructor{Some}} \coqdocvar{m}, então \coqdocvar{m} pertence a \coqdocvar{l}. Isso é
    necessário para garantir que \coqdocdefinition{remove\_one} de fato diminui o tamanho da
    lista (usado na medida da função \coqdoclemma{ss} mais abaixo). 
\begin{coqdoccode}
\coqdocnoindent
\coqdockw{Lemma} \coqdocvar{select\_min\_in} : \coqdockw{\ensuremath{\forall}} \coqdocvar{l} \coqdocvar{m}, \coqdocvar{select\_min} \coqdocvar{l} = \coqdocvar{Some} \coqdocvar{m} \ensuremath{\rightarrow} \coqdocvar{In} \coqdocvar{m} \coqdocvar{l}.\coqdoceol
\coqdocemptyline
\end{coqdoccode}
Corolário: se \coqdoclemma{select\_min} \coqdocvar{l} = \coqexternalref{None}{http://rocq-prover.org/doc/V9.2/corelib//Corelib.Init.Datatypes}{\coqdocconstructor{None}}, então \coqdocvar{l} é vazia (o único caso da
    definição que retorna None é o caso nil). 
\begin{coqdoccode}
\coqdocnoindent
\coqdockw{Lemma} \coqdocvar{select\_min\_none} : \coqdockw{\ensuremath{\forall}} \coqdocvar{l}, \coqdocvar{select\_min} \coqdocvar{l} = \coqdocvar{None} \ensuremath{\rightarrow} \coqdocvar{l} = \coqdocvar{nil}.\coqdoceol
\coqdocemptyline
\coqdocemptyline
\end{coqdoccode}
\coqdocdefinition{remove\_one} \coqdocvar{x} \coqdocvar{l} remove UMA ocorrência do elemento \coqdocvar{x} da lista \coqdocvar{l}.
    Esta função corrige o bug da definição original de \coqdoclemma{ss}, que recursava
    sobre a cauda \coqdocvar{tl} em vez de remover o mínimo efetivamente encontrado. 
\begin{coqdoccode}
\coqdocnoindent
\coqdockw{Fixpoint} \coqdocvar{remove\_one} (\coqdocvar{x} : \coqdocvar{nat}) (\coqdocvar{l} : \coqdocvar{list} \coqdocvar{nat}) : \coqdocvar{list} \coqdocvar{nat} :=\coqdoceol
\coqdocindent{1.00em}
\coqdockw{match} \coqdocvar{l} \coqdockw{with}\coqdoceol
\coqdocindent{1.00em}
\ensuremath{|} \coqdocvar{nil} \ensuremath{\Rightarrow} \coqdocvar{nil}\coqdoceol
\coqdocindent{1.00em}
\ensuremath{|} \coqdocvar{h} :: \coqdocvar{t} \ensuremath{\Rightarrow} \coqdockw{if} \coqdocvar{h} =? \coqdocvar{x} \coqdockw{then} \coqdocvar{t} \coqdockw{else} \coqdocvar{h} :: \coqdocvar{remove\_one} \coqdocvar{x} \coqdocvar{t}\coqdoceol
\coqdocindent{1.00em}
\coqdockw{end}.\coqdoceol
\coqdocemptyline
\coqdocnoindent
\coqdockw{Lemma} \coqdocvar{remove\_one\_length} : \coqdockw{\ensuremath{\forall}} \coqdocvar{x} \coqdocvar{l}, \coqdocvar{In} \coqdocvar{x} \coqdocvar{l} \ensuremath{\rightarrow} \coqdocvar{length} (\coqdocvar{remove\_one} \coqdocvar{x} \coqdocvar{l}) < \coqdocvar{length} \coqdocvar{l}.\coqdoceol
\coqdocemptyline
\coqdocnoindent
\coqdockw{Lemma} \coqdocvar{remove\_one\_In} : \coqdockw{\ensuremath{\forall}} \coqdocvar{y} \coqdocvar{x} \coqdocvar{l}, \coqdocvar{In} \coqdocvar{y} (\coqdocvar{remove\_one} \coqdocvar{x} \coqdocvar{l}) \ensuremath{\rightarrow} \coqdocvar{In} \coqdocvar{y} \coqdocvar{l}.\coqdoceol
\coqdocemptyline
\coqdocnoindent
\coqdockw{Lemma} \coqdocvar{remove\_one\_perm} : \coqdockw{\ensuremath{\forall}} \coqdocvar{x} \coqdocvar{l}, \coqdocvar{In} \coqdocvar{x} \coqdocvar{l} \ensuremath{\rightarrow} \coqdocvar{Permutation} \coqdocvar{l} (\coqdocvar{x} :: \coqdocvar{remove\_one} \coqdocvar{x} \coqdocvar{l}).\coqdoceol
\coqdocemptyline
\coqdocemptyline
\end{coqdoccode}
Se \coqdocvar{l} está ordenada e \coqdocvar{x} é menor ou igual a todos os elementos de \coqdocvar{l},
    então \coqdocvar{x}::\coqdocvar{l} também está ordenada. 
\begin{coqdoccode}
\coqdocnoindent
\coqdockw{Lemma} \coqdocvar{le\_all\_sorted} : \coqdockw{\ensuremath{\forall}} \coqdocvar{l} \coqdocvar{x}, \coqdocvar{Sorted} \coqdocvar{le} \coqdocvar{l} \ensuremath{\rightarrow} \coqdocvar{le\_all} \coqdocvar{x} \coqdocvar{l} \ensuremath{\rightarrow} \coqdocvar{Sorted} \coqdocvar{le} (\coqdocvar{x} :: \coqdocvar{l}).\coqdoceol
\coqdocemptyline
\coqdocemptyline
\end{coqdoccode}
ATENÇÃO / correção do enunciado: a versão original do professor definia
        ss (h::tl) := match select\_min l with Some m => m :: ss tl end
    ou seja, recursava sobre \coqdocvar{tl} (a lista SEM a cabeça), e não sobre a lista
    sem o MÍNIMO encontrado. Isso quebra a permutação: por exemplo,
    ss 4;2;1;3 geraria 1;1;1;3 em vez de 1;2;3;4.
    A correção é recursar sobre \coqdocdefinition{remove\_one} \coqdocvar{m} \coqdocvar{l}, removendo de fato o
    elemento mínimo encontrado. 
\begin{coqdoccode}
\coqdocnoindent
\coqdockw{Function} \coqdocvar{ss} (\coqdocvar{l} : \coqdocvar{list} \coqdocvar{nat}) \{\coqdockw{measure} \coqdocvar{length} \coqdocvar{l}\} : \coqdocvar{list} \coqdocvar{nat} :=\coqdoceol
\coqdocindent{1.00em}
\coqdockw{match} \coqdocvar{select\_min} \coqdocvar{l} \coqdockw{with}\coqdoceol
\coqdocindent{1.00em}
\ensuremath{|} \coqdocvar{None} \ensuremath{\Rightarrow} \coqdocvar{nil}\coqdoceol
\coqdocindent{1.00em}
\ensuremath{|} \coqdocvar{Some} \coqdocvar{m} \ensuremath{\Rightarrow} \coqdocvar{m} :: \coqdocvar{ss} (\coqdocvar{remove\_one} \coqdocvar{m} \coqdocvar{l})\coqdoceol
\coqdocindent{1.00em}
\coqdockw{end}.\coqdoceol
\coqdocemptyline
\end{coqdoccode}
Teorema final: \coqdoclemma{ss} retorna uma permutação ordenada da lista de entrada. 
\begin{coqdoccode}
\coqdocnoindent
\coqdockw{Theorem} \coqdocvar{selectionsort\_correct} : \coqdockw{\ensuremath{\forall}} \coqdocvar{l}, \coqdocvar{Sorted} \coqdocvar{le} (\coqdocvar{ss} \coqdocvar{l}) \ensuremath{\land} \coqdocvar{Permutation} \coqdocvar{l} (\coqdocvar{ss} \coqdocvar{l}).\coqdoceol
\coqdocemptyline
\end{coqdoccode}
testes: 
\begin{coqdoccode}
\coqdocemptyline
\coqdocnoindent
\coqdockw{Eval} \coqdoctac{compute} \coqdoctac{in} \coqdocvar{ss} [4;2;1;3].\coqdoceol
\coqdocnoindent
\coqdockw{Eval} \coqdoctac{compute} \coqdoctac{in} \coqdocvar{ss} [5;4;3;2;1].\coqdoceol
\coqdocnoindent
\coqdockw{Eval} \coqdoctac{compute} \coqdoctac{in} \coqdocvar{ss} (\coqdocvar{nil} : \coqdocvar{list} \coqdocvar{nat}).\coqdoceol
\coqdocemptyline
\coqdocnoindent
\coqdockw{Lemma} \coqdocvar{teste1} : \coqdocvar{ss} [4;2;1;3] = [1;2;3;4].\coqdoceol
 \coqdocemptyline
\coqdocnoindent
\coqdockw{Lemma} \coqdocvar{teste2} : \coqdocvar{ss} [5;4;3;2;1] = [1;2;3;4;5].\coqdoceol
 \coqdocemptyline
\coqdocnoindent
\coqdockw{Lemma} \coqdocvar{teste3} : \coqdocvar{ss} (@\coqdocvar{nil} \coqdocvar{nat}) = \coqdocvar{nil}.\coqdoceol
 \end{coqdoccode}


\end{document}